\chapter*{이 책을 읽는 방법}
\addcontentsline{toc}{chapter}{이 책을 읽는 방법}
\markboth{이 책을 읽는 방법}{이 책을 읽는 방법}

Volume 3은 하나의 긴 증명만을 향해 직선으로 진행하지 않는다. 다항식의 갈루아군을 실제로 계산하는 흐름, 원분체와 체확장의 불변량을 익히는 흐름, 근호해법의 가능성과 한계를 밝히는 흐름, 고전 작도문제를 해결하는 흐름이 서로 교차한다. 이 절은 장 사이의 의존관계와 반복되는 개념의 역할을 먼저 보여 준다.

\section*{선수개념}

다음 내용은 Volume 2에서 준비되어 있다고 가정한다.
\begin{itemize}
  \item 분해체, 대수적 폐포와 최소다항식
  \item 분리확장과 정규확장
  \item 유한 갈루아 확장과 갈루아군
  \item 부분군과 중간체의 대응, 정규부분군과 몫군
  \item 유한체와 프로베니우스 자동동형
\end{itemize}
필요한 군론은 Volume 1의 부분군, 정규부분군, 몫군, 군작용과 궤도--안정자 공식이다. 도함열과 가해군처럼 새로 필요한 군론은 Chapter 6에서 다시 정의한다.

\section*{장 사이의 의존관계}

전권의 주된 흐름은 다음과 같다.
\[
\begin{array}{ccccccc}
0&\longrightarrow&1&\longrightarrow&2&\longrightarrow&3\\
&&\downarrow&&&&\downarrow\\
&&4&\longrightarrow&5&\longrightarrow&6\longrightarrow7\\
&&\downarrow\\
&&8
\end{array}
\]
이 그림은 엄밀한 선수관계를 모두 표시한 것이 아니라 학습의 중심 경로를 요약한다. Chapter 8은 특히 Chapter 1의 갈루아군과 Chapter 4의 원분체를 사용한다.

\begin{description}[style=nextline,leftmargin=0pt]
  \item[Chapter 0--1: 방정식의 대칭을 읽는 언어]
  분해체의 자동동형을 근의 순열로 바꾸고, 기약성과 추이성을 연결한다. 뒤의 모든 계산에서 출발점이 된다.

  \item[Chapter 2--3: 계수에서 갈루아군을 추적하는 도구]
  대칭다항식, 판별식, 종결식과 삼차분해식을 이용하여 사차와 오차다항식의 갈루아군을 계산한다.

  \item[Chapter 4--5: 순환확장과 체의 불변량]
  원분체, 힐베르트 정리 90, 노름, 대각합과 체매장의 선형독립을 전개한다. Chapter 6의 근호확장 구조를 위한 재료이다.

  \item[Chapter 6--7: 근호해법의 가능성과 한계]
  가해군과 근호탑의 동치를 세운 뒤, 일반다항식의 갈루아군 $S_n$과 $A_n$의 단순성을 결합하여 아벨--루피니 정리를 증명한다.

  \item[Chapter 8: 갈루아 이론의 기하학적 응용]
  이차확장 탑과 갈루아 $2$-군을 자와 컴퍼스 작도로 번역한다. Chapter 4의 원분체가 정다각형 작도 판정에 다시 등장한다.
\end{description}

\section*{목적에 따른 네 가지 학습 경로}

\begin{description}[style=nextline,leftmargin=0pt]
  \item[전권 정독]
  Chapter 0부터 8까지 순서대로 읽는다. 정의와 정리의 의존관계가 가장 자연스럽게 이어진다.

  \item[갈루아군 계산 중심]
  Chapter 0, 1, 2, 3을 먼저 읽고 Chapter 6의 Cardano 공식과 사차공식 부분으로 간다. 수치계수 다항식의 기약성, 판별식, 법 $p$ 인수분해와 분해식을 이용하는 절차에 초점을 둔다.

  \item[근호해법과 아벨--루피니 중심]
  Chapter 1의 순열표현을 복습한 뒤 Chapter 4, 5, 6, 7을 읽는다. 순환확장, 노름과 대각합, 가해군, 일반다항식의 순서가 핵심이다.

  \item[정다각형과 작도 중심]
  Chapter 1의 분해체와 갈루아군, Chapter 4의 원분체를 읽고 Chapter 8로 간다. 작도가능성의 완전한 판정에는 최소다항식의 차수뿐 아니라 정상폐포의 갈루아군이 필요하다.
\end{description}

\section*{반복되는 내용은 왜 다시 나오는가}

몇몇 정리는 서로 다른 목적 때문에 두 번 등장한다.
\begin{itemize}
  \item 일반다항식의 갈루아군 $S_n$은 Chapter 1에서 첫 계산 모형으로 소개하고, Chapter 7에서 아벨--루피니 정리의 정확한 입력으로 다시 정식화한다.
  \item 삼차분해식은 Chapter 3에서 사차 갈루아군을 판정하는 도구이고, Chapter 6에서는 사차근을 근호로 복원하는 도구이다.
  \item 체매장의 선형독립은 Chapter 4에서 힐베르트 정리 90에 필요한 형태로 먼저 증명하고, Chapter 5에서 캐릭터 선형독립의 일반 정리로 확장한다.
\end{itemize}
두 번째 등장은 단순한 반복이 아니라 같은 대상을 다른 기능으로 사용하는 단계이다.

\section*{연습문제 체계}

Chapter 2--8에는 각 장마다 30개의 문제가 있고, 입문 장인 Chapter 1에는 종합문제 하나를 더 둔 31개의 문제가 있다.
\begin{description}[style=nextline,leftmargin=0pt]
  \item[A. 계산과 기본 판정]
  정의와 표준 알고리즘을 실제 다항식과 체확장에 적용한다.
  \item[B. 핵심 정리의 증명]
  본문 정리의 논증을 독자가 다시 구성하거나 빠진 단계를 완성한다.
  \item[C. 구조와 종합]
  여러 장의 결과를 결합하거나 새로운 예제에 이론을 적용한다.
  \item[별표 문제]
  뒤의 선별 풀이 절에 상세 풀이가 있다.
\end{description}
전권에는 연습문제 241개와 선별 상세 풀이 121개가 있다. 처음 읽을 때에는 각 절의 이해 점검과 A 문제를 먼저 풀고, 두 번째 읽기에서 B와 C 문제로 넘어가는 편이 좋다.

\section*{기호와 용어의 약속}

\begin{itemize}
  \item $D_n$은 정$n$각형의 대칭군, 즉 위수 $2n$인 이면체군을 뜻한다. 따라서 $D_4$의 위수는 $8$, $D_5$의 위수는 $10$이다.
  \item 사차식에 붙이는 resolvent cubic은 전권에서 \emph{삼차분해식}이라 부른다.
  \item $\Nm_{L/K}$는 노름, $\Tr_{L/K}$는 대각합이다.
  \item normal extension은 \emph{정규확장}, normal closure는 이 시리즈의 기존 용어에 따라 \emph{정상폐포}라 부른다.
  \item ``일반 $n$차다항식''은 계수들이 대수적으로 독립인 보편 다항식을 뜻한다. 모든 수치계수 다항식의 갈루아군이 $S_n$이라는 뜻이 아니다.
  \item 아벨--루피니 정리는 모든 오차방정식이 근호로 풀리지 않는다고 말하지 않는다. 보편적인 오차 근호공식이 없다는 정리이며, 개별 오차식은 갈루아군이 가해군이면 근호로 풀린다.
\end{itemize}

\begin{keyidea}
이 권을 관통하는 질문은 하나이다.
\[
  \boxed{\text{계수로부터 보이는 정보가 근들의 가능한 순열을 얼마나 제한하는가?}}
\]
판별식, 분해식, 프로베니우스 순환형, 노름과 대각합, 도함열과 작도가능성은 모두 이 질문에 서로 다른 방식으로 답한다.
\end{keyidea}
